\documentclass[11pt]{article}

\usepackage[margin=1in]{geometry}
\usepackage{amsmath,amssymb,amsfonts,amsthm}
\usepackage{array}
\usepackage{booktabs}
\usepackage{graphicx}
\usepackage{hyperref}
\usepackage{url}
\usepackage{xcolor}
\emergencystretch=4em

\hypersetup{
  colorlinks=true,
  linkcolor=blue,
  citecolor=blue,
  urlcolor=blue
}

\newtheorem{assumption}{Assumption}
\newtheorem{lemma}{Lemma}
\newtheorem{theorem}{Theorem}
\newtheorem{corollary}{Corollary}

\newcommand{\method}{Gauss6\slash FullVA}
\newcommand{\tfe}{TFE}
\newcommand{\artifact}[1]{\nolinkurl{#1}}
\newcommand{\figpath}{figures}
\newcolumntype{P}[1]{>{\raggedright\arraybackslash}p{#1}}
\newcommand{\paperfig}[3]{%
  \begin{figure}[htbp]
    \centering
    \IfFileExists{#1}{%
      \includegraphics[width=0.92\linewidth]{#1}%
    }{%
      \fbox{\parbox{0.88\linewidth}{Missing figure: \nolinkurl{#1}}}%
    }
    \caption{#2}
    \label{#3}
  \end{figure}
}

\title{A Conditional Sixth-Order Lie-Group FullVA Integrator for Lower-Pair Mechanisms}
\author{Jingquan Wang}
\date{v047 concise artifact-backed draft}

\begin{document}
\maketitle

\begin{abstract}
This concise draft states only the accepted v047 paper claim.  The production
integrator is the \method{} one-step map, a conditional sixth-order Gauss
collocation FullVA method on a Lie-group lower-pair formulation.  On the smooth
cylindrical-chain benchmark it records observed position and velocity orders
7.161 and 7.066.  The same artifact contract accepts four ASME-style examples:
single pendulum, double pendulum, four link, and slider crank.  The local
paper-style $m=3$ Gauss--Lobatto temporal finite-element target used as the
paper-facing comparator has expected order five.  Therefore the accepted
result is a conditional formal-order comparison: the accepted \method{} path
has a conditional order-six theorem on the selected smooth branch, the local
paper-style target is order five, and the four examples provide separated
method-side evidence roles rather than a single undifferentiated validation
claim.

The draft deliberately does not claim a complete reproduction of the source
paper's \tfe{} nonlinear residual.  In this repository, independent full-\tfe{}
replacement means replacing all 132 Gauss/FullVA stage rows by paper-derived
\tfe{} weak rows inside Newton, without endpoint source data, terminal-row
replacement, projection, or a target-direction oracle.  That stronger gate
remains open as \nolinkurl{full_tfe_stage_replacement=false}.
\end{abstract}

\section{Accepted Claim}

The v047 evidence supports a narrow and useful statement: \method{} is a
conditional sixth-order method compared with the local paper-style $m=3$
Gauss--Lobatto \tfe{} formula target under the current artifact scope.  The
claim is not that every row of the source paper's full \tfe{} residual has been reproduced in the
nonlinear solver.  This distinction keeps the paper aligned with the actual
validation evidence.

\begin{table}[htbp]
\centering
\caption{Concise claim boundary.}
\label{tab:claim-boundary}
\begin{tabular}{P{0.27\linewidth}P{0.59\linewidth}}
\toprule
Item & Current status \\
\midrule
Accepted method &
\method{}, a conditional sixth-order Gauss collocation/FullVA one-step map on
the selected smooth branch. \\
Observed smooth order &
Position/velocity orders 7.161/7.066 over
$h=\{0.04,0.02,0.01\}$ against $h=0.005$. \\
Comparator &
Local paper-style $m=3$ Gauss--Lobatto \tfe{} formula target with expected
order five. \\
Four-example gate &
Single and double pendula provide dynamic-order evidence; four link and slider
crank provide lower-pair mechanism-coverage diagnostics. \\
Open reproduction gate &
\nolinkurl{full_tfe_stage_replacement=false}. \\
\bottomrule
\end{tabular}
\end{table}

\section{Integrator}

Let $x_n=(q_n,v_n,a_n)$ denote the endpoint position, velocity, and
acceleration variables, with orientations represented on the local Lie-group
chart used by the v047 solver.  The accepted one-step map is
\begin{equation}
  x_{n+1}=\Psi_h^{\mathrm{G6FVA}}(x_n),
\end{equation}
where the stage variables $Z$ solve the 132-row nonlinear system
\begin{equation}
  R_{\mathrm{G6FVA}}(q_n,v_n,a_n,Z;h)=0.
\end{equation}
The lower-pair constraints enter at position, velocity, and acceleration
levels,
\begin{equation}
  \Phi(q)=0,\qquad
  \Phi_q(q)v=0,\qquad
  \Phi_q(q)a+\dot{\Phi}_q(q,v)v=0.
\end{equation}
The endpoint is reconstructed by Gauss collocation quadrature and the local
Lie-group retraction.  Endpoint velocity closure is audited by the generated
KKT/projection artifacts.  This closure is part of the accepted Gauss/FullVA
evidence path; it is not an accepted independent full-\tfe{} replacement.

\begin{table}[htbp]
\centering
\caption{Operational definition of the accepted one-step map.}
\label{tab:algorithm}
\begin{tabular}{P{0.16\linewidth}P{0.70\linewidth}}
\toprule
Step & Operation \\
\midrule
1 &
Assemble cylindrical lower-pair constraints and FullVA stage variables for
state $x_n$ and step size $h$. \\
2 &
Solve the square 132-row \method{} residual to the Newton tolerance recorded in
the CSV/JSON artifacts. \\
3 &
Recover the endpoint by the Gauss quadrature and Lie-group retraction used by
the v047 trajectory code. \\
4 &
Audit endpoint constraint and velocity closure with the KKT/projection
diagnostics. \\
5 &
Record residuals, constraints, observed orders, figures, and validator-visible
summary fields. \\
\bottomrule
\end{tabular}
\end{table}

\section{Proof Sketch}

\begin{assumption}[Accepted-path regularity]
\label{ass:regularity}
On the selected smooth branch, the compact-tube regularity, endpoint map,
implementation binding, and scaled nonlinear residual condition
\(\eta_h^{\rm tube}\le c_\eta h^7\) hold.  The implemented 132-row FullVA
residual bridge is the direct route used by the strict proof: the 96
non-dynamic rows supply the row-local \(O(h^7)\) input, and the 36
Newton--Euler rows close by direct substitution on the lifted Gauss stage.
\end{assumption}

\begin{theorem}[Order of the accepted \method{} map]
\label{thm:g6fullva-order}
Under Assumption~\ref{ass:regularity}, read as the retained accepted-branch
compact-tube/P6 solver-envelope interface together with the P4/P5 direct
\(96+36\) residual bridge, the accepted \method{} one-step map has local
defect \(O(h^7)\) and reduced-chart/reported position--velocity grid error
\(O(h^6)\) on fixed smooth time intervals for reported grids satisfying the
predeclared branch, endpoint, proof-norm, and
\(\eta_h^{\rm tube}\le c_\eta h^7\) conditions.  For driven rows, the
time-indexed residual, Gauss stage times, endpoint map, and proof norm are
fixed on the same transition slice before constants or diagnostics are used.
This is a conditional
order-six statement for the accepted branch-selected map only; it retains P6
as a theorem-level solver-envelope condition, does not prove a P7
residual-to-error transfer, and does not use source-policy rows or the optional
primitive/Taylor route.  Equivalently, the theorem asserts only the local defect
$O(h^7)$ and same-initial-state global error $O(h^6)$ under these retained
interfaces and only for this reduced/reporting grid scope.
\end{theorem}

\begin{proof}
The strict CMAME proof route is not a generic inheritance statement.  The
classical three-stage Gauss local defect is inserted into the implemented
FullVA residual on the selected branch.  The direct 132-row bridge gives an
\(O(h^7)\) residual defect on the lifted Gauss stage; endpoint closure and the
scaled inexact-Newton residual add \(O(h^7)\) perturbations under the retained
compact-tube hypotheses.  The standard local-to-global transfer for the
accepted branch then gives the \(O(h^6)\) grid estimate on the same reported
grid.  The proof follows the constrained Lie-group BDF proof order only as an
ordering discipline: the discrete object is fixed first, the smooth branch is
inserted next, the local defect is proved for that object, and the global
recurrence is used only afterward.  No BDF local-error, hidden-constraint,
multiplier-history, or multistep stability estimate is imported.  The solver
bound remains a retained P6 envelope fixed before the run,
and the P7 residual-to-error interface remains open.  Residual-only mechanism
rows, source-policy rows, and the separate primitive/Taylor route are not used
as theorem inputs and cannot be spliced into the direct bridge.
\end{proof}

\begin{lemma}[Order of the local paper-style \tfe{} target]
\label{lem:tfe-target-order}
The proof-level order target for the $m=3$ Gauss--Lobatto paper-style
\tfe{} formula used in this manuscript is five.
\end{lemma}

\begin{proof}
The source-paper family is compared through the temporal finite-element
formula-order convention $p_{\mathrm{TFE}}(m)=2m-1$.  Substituting $m=3$
gives $p_{\mathrm{TFE}}(3)=5$.  This is distinct from the source paper's
reported DAE pendulum order reduction.
\end{proof}

\begin{theorem}[Accepted formal-order comparison theorem]
\label{thm:formal-order}
Under Assumption~\ref{ass:regularity} and the current v047 artifact contract,
the \method{} path is accepted as a conditional sixth-order method relative
to the local paper-style $m=3$ Gauss--Lobatto \tfe{} formula target with
expected order five.
\end{theorem}

\begin{proof}
Theorem~\ref{thm:g6fullva-order} gives the conditional order-six method claim
for the accepted \method{} path through the direct 132-row residual bridge.
Lemma~\ref{lem:tfe-target-order} gives the fifth-order local paper-style
comparator.  The generated h-sweep records
observed smooth position/velocity orders 7.161/7.066, which are above that
target.  The read-only four-example validator accepts single pendulum, double
pendulum, four link, and slider crank on the same method path.  Therefore the
formal-order comparison is accepted inside the artifact scope.  The theorem
excludes complete source-paper residual reproduction because the
machine-readable boundary still records
\texttt{full\_tfe\_stage\_\allowbreak replacement=false}.
% validator token: full_tfe_stage_replacement=false
\end{proof}

\begin{corollary}[What is not proved]
The theorem does not prove that the implemented nonlinear stage residual is the
complete source-paper \tfe{} residual.
\end{corollary}

\begin{proof}
The accepted method path solves $R_{\mathrm{G6FVA}}=0$.  A full source-paper
\tfe{} reproduction claim would require an accepted residual
$R_{\tfe{}}=0$ using paper-derived weak rows for all 132 stage rows inside
Newton, with the same rank, residual, endpoint, and h-sweep checks.  That gate
is explicitly open.
\end{proof}

\begin{table}[htbp]
\centering
\caption{Proof obligations and artifact evidence for
Theorem~\ref{thm:formal-order}.}
\label{tab:proof-evidence}
\begin{tabular}{P{0.29\linewidth}P{0.53\linewidth}}
\toprule
Proof obligation & Current artifact evidence \\
\midrule
Accepted one-step map &
\artifact{CLAIM_BOUNDARY.json} and \artifact{summary_v047.json} record the
\method{} 132-row stage-residual boundary. \\
Sixth-order smooth behavior &
\artifact{cylindrical_chain_convergence.csv} records smooth
position/velocity orders 7.161/7.066 over $h=[0.04,0.02,0.01]$ against
\artifact{reference_h=0.005}. \\
Comparator order &
\artifact{SOURCE_PAPER_COMPARISON.md} records the local $m=3$
Gauss--Lobatto \tfe{} expected order $2m-1=5$. \\
Four-example support &
\artifact{validate_four_asme_minimal.py} accepts single pendulum, double
pendulum, four link, and slider crank. \\
Non-claim boundary &
\artifact{PROOF_EVIDENCE_MATRIX.md} and the submission validators preserve
\artifact{full_tfe_stage_replacement=false}. \\
\bottomrule
\end{tabular}
\end{table}

\section{Numerical Evidence}

The generated convergence artifacts use $h=\{0.04,0.02,0.01\}$ and an
$h=0.005$ reference.  Table~\ref{tab:orders} reports the headline order
evidence.

\begin{table}[htbp]
\centering
\caption{Observed orders for the accepted \method{} path.}
\label{tab:orders}
\begin{tabular}{P{0.25\linewidth}P{0.22\linewidth}P{0.20\linewidth}P{0.20\linewidth}}
\toprule
Case & Step sizes & Position order & Velocity order \\
\midrule
Smooth cylindrical chain & $0.04,0.02,0.01$ & 7.161 & 7.066 \\
Sharp Brown--McPhee case & $0.04,0.02,0.01$ & 1.894 & 2.683 \\
\bottomrule
\end{tabular}
\end{table}

\paperfig{\figpath/convergence.png}
{Convergence of the accepted \method{} path.  The smooth branch is the
paper-facing high-order evidence; the sharp branch is reported as a practical
coarse-regime caveat.}
{fig:convergence}

The sharp case does not contradict the smooth formal-order comparison.  It is
reported as a practical caveat: coarse sharp-friction steps are order-reduced,
while deeper fixed refinement recovers high order at increased cost.

\section{Four ASME-Style Examples}

The same method-side gate accepts the four examples listed in
Table~\ref{tab:asme}.  The single and double pendulum rows provide order-based
method evidence.  The four-link and slider-crank rows provide closed-loop
kinematic FullVA solves and reaction-dynamics checks.

\begin{table}[htbp]
\centering
\caption{Four-example evidence summary.}
\label{tab:asme}
\begin{tabular}{P{0.22\linewidth}P{0.27\linewidth}P{0.37\linewidth}}
\toprule
Example & Accepted evidence & Headline value \\
\midrule
Single pendulum &
Driven absolute-coordinate FullVA residual &
Minimum observed order 6.024; max stage residual $3.664\times10^{-12}$. \\
Double pendulum &
Method-side FullVA reference policy &
Minimum orientation/angular-velocity order 6.089; endpoint constraints below
$8.323\times10^{-12}$. \\
Four link &
Closed-loop kinematic FullVA and reaction dynamics &
Max constraint norm $1.052\times10^{-14}$; max dynamics residual
$1.338\times10^{-13}$. \\
Slider crank &
Closed-loop kinematic FullVA and reaction dynamics &
Max constraint norm $1.204\times10^{-15}$; max dynamics residual
$6.492\times10^{-15}$. \\
\bottomrule
\end{tabular}
\end{table}

\paperfig{\figpath/asme_lower_pair_graph_bridge.png}
{Lower-pair graph bridge used to keep the four ASME-style examples under the
same validation vocabulary.}
{fig:asme-graph}

\paperfig{\figpath/asme_closed_loop_kinematic_fullva.png}
{Closed-loop kinematic FullVA evidence for the four-link and slider-crank
examples.}
{fig:closed-loop}

\section{Comparator and Boundary}

The comparator is intentionally local: the encoded $m=3$ Gauss--Lobatto \tfe{}
formula target has expected order five.  It is sufficient for the paper-facing
formal-order comparison because the accepted method has a conditional
order-six theorem and the evidence package includes four separated example
roles.  It is not the same as an accepted
implementation of every source-paper residual row.

The source-paper order convention used by the local comparator is
\begin{equation}
  p_{\mathrm{TFE}}(m)=2m-1.
\end{equation}
Thus the local $m=3$ Gauss--Lobatto \tfe{} target used here has
$p_{\mathrm{TFE}}(3)=5$.  The accepted \method{} path has method-order claim
six and observed smooth orders 7.161/7.066.  The positive comparison is
therefore order-six \method{} versus an order-five paper-style target.  The
comparison does not assert that the source paper contains this repository's
132-row \method{} residual, nor that this repository has accepted a full
source-paper \tfe{} residual substitution.  The standalone comparison note
\artifact{SOURCE_PAPER_COMPARISON.md} records this boundary.

The example-level order boundary is recorded separately in
\artifact{ORDER_ACCEPTANCE_GATE.md}.  It keeps \artifact{single_pendulum} and
\artifact{double_pendulum} as accepted dynamic order rows and keeps
\artifact{four_link} and \artifact{slider_crank} as mechanism-coverage rows
that are not accepted external dynamic-order rows.  The same gate records the
\artifact{coarse_first_no_default_1e-4} execution policy.

\begin{table}[htbp]
\centering
\caption{Comparator versus open reproduction gate.}
\label{tab:comparator}
\begin{tabular}{P{0.27\linewidth}P{0.57\linewidth}}
\toprule
Object & Role \\
\midrule
Local $m=3$ Lobatto-\tfe{} formula target &
Order-five comparator for the formal-order comparison. \\
Accepted \method{} path &
Conditional order-six method path used for the paper's positive claim. \\
Independent full-\tfe{} replacement &
Open stronger reproduction target requiring paper-derived rows inside Newton
for all 132 stage rows. \\
\bottomrule
\end{tabular}
\end{table}

\section{Limitations}

Three caveats remain explicit in the claim boundary.
\begin{itemize}
  \item Sparse AD structure is correct and quantified, but dense
  \artifact{jacfwd} is still faster in the current wall-clock evidence.
  \item Independent full-\tfe{} stage replacement is not accepted:
  \nolinkurl{full_tfe_stage_replacement=false}.
  \item The sharp-friction coarse regime is order-reduced, although deeper
  fixed refinement recovers high order at practical cost.
\end{itemize}
These limitations do not invalidate the narrower formal-order comparison.
They define future engineering and reproduction work.

\section{Reproducibility Contract}

The concise paper is checked from existing artifacts.  Routine validation does
not require the full \artifact{run_v047.py} generator.  The relevant read-only
commands are:
\begin{verbatim}
../.venv_sbel/bin/python validate_concise_paper.py
../.venv_sbel/bin/python validate_source_paper_comparison.py
../.venv_sbel/bin/python validate_order_acceptance_gate.py
../.venv_sbel/bin/python validate_paper_package.py
../.venv_sbel/bin/python ../validate_pipeline_outputs.py
\end{verbatim}
The machine-readable claim boundary is \artifact{CLAIM_BOUNDARY.json}; the
top-level current-state boundary is \artifact{CURRENT_PIPELINE_CONTRACT.md}.

\section{Conclusion}

The accepted v047 result is simple when the diagnostic replacement work is
kept separate from the method claim.  The \method{} path follows the
conditional sixth-order Lie-group FullVA theorem, records smooth observed
orders 7.161/7.066, and provides four-example method-side coverage.  Against
the local order-five paper-style \tfe{} comparator, this supports the
formal-order comparison.  Complete source-paper \tfe{} residual replacement
remains future work.

\end{document}
